% id: 85-20
% book: 베이직쎈 미적분1 · 05 도함수의 활용 (2)
% section: 개념 38 함수의 최대와 최소
% passage: 주어진 구간에서 다음 함수 $f(x)$의 최댓값과 최솟값을 구하시오.
% figure: tikz:fig-85-20
% answer: 
% answer_source: 
% variant_level: 0 (원본 전사)
% 이 파일은 build.mjs 가 items.json 에서 만든다. 고칠 때는 items.json 을 고친다.
\smash{\raisebox{1.5mm}{\dmkichul{베이직쎈 미적분1 85쪽 20번}}}
\noindent\begin{minipage}[t]{0.58\linewidth}\vspace{0pt}\raggedright $f(x)=2x^3-6x^2+4$ \qquad $[-2{,}\mkern3mu\mathopen{}2]$\end{minipage}\hfill\begin{minipage}[t]{0.4\linewidth}\vspace{0pt}\centering \resizebox{\ifdim\width>\linewidth\linewidth\else\width\fi}{!}{% 85-20(85쪽) — 20번 풀이 힌트 오른쪽의 그림 · f(x)=2x^3-6x^2+4 의 그래프(닫힌구간 [-2,2] 밖은 점선). 스캔 화질이 낮아 TikZ 로 다시 그림.
% 원문에 있는 것만: 라벨 4, -4, -36, -2, 2, O, x, y, y=f(x) 와 안내 점선. 눈금·원문에 없는 값은 만들지 않는다.
\begin{tikzpicture}[x=0.55cm, y=0.105cm, line width=0.6pt]
  \draw[->, >=stealth, line width=0.7pt] (-2.8,0) -- (3.2,0) node[below=1pt, font=\small] {$x$};
  \draw[->, >=stealth, line width=0.7pt] (0,-45) -- (0,11) node[left=1pt, font=\small] {$y$};
  \draw[densely dotted, line width=0.4pt] (-2,0) -- (-2,-36);
  \draw[densely dashed, line width=0.4pt] (-2,-36) -- (0,-36);
  \draw[densely dashed, line width=0.4pt] (0.35,-4) -- (2,-4);
  \draw[line width=0.4pt] (2,0) -- (2,-4);
  \draw[densely dashed, line width=0.6pt, domain=-2.2:-2, samples=20, smooth] plot (\x, {2*\x*\x*\x-6*\x*\x+4});
  \draw[line width=0.6pt, domain=-2:2, samples=90, smooth] plot (\x, {2*\x*\x*\x-6*\x*\x+4});
  \draw[densely dashed, line width=0.6pt, domain=2:2.9, samples=30, smooth] plot (\x, {2*\x*\x*\x-6*\x*\x+4});
  \node[left=1pt, font=\small] at (0,4.6) {$4$};
  \node[above=1pt, font=\small] at (-2,0) {$-2$};
  \node[above=1pt, font=\small] at (2,0) {$2$};
  \node[below right=1pt and 1pt, font=\small] at (0,-4) {$-4$};
  \node[right=2pt, font=\small] at (0,-36) {$-36$};
  \node[below left=-1pt and -2pt, font=\small] at (0,0) {$\pt{O}$};
  \node[font=\small] at (1.8,9.5) {$y=f(x)$};
\end{tikzpicture}
}\end{minipage}\par
\dmhint{$f(x)=2x^3-6x^2+4$에서\\ \hspace*{1.5em}$f'(x)=6x^2-12x=6x(x-2)$\\ $f'(x)=0$에서 \qquad $x=\blank$ 또는 $x=2$\par\smallskip\begin{center}\begin{tabular}{c|c|c|c|c|c}\hline $x$ & $-2$ & $\cdots$ & \blank & $\cdots$ & $2$ \\ \hline $f'(x)$ & & $+$ & $0$ & $-$ & $0$ \\ \hline $f(x)$ & \blank[2.4em] & $\nearrow$ & \blank & $\searrow$ & $-4$ \\ \hline\end{tabular}\end{center}\smallskip\noindent 따라서 함수 $f(x)$는 $x=\blank$에서 최댓값 \blank, $x=\blank$에서 최솟값 \blank을 갖는다.}
